\section{Concluding Remarks}

\subsection{Analytic Formulations and Arithmetic Consequences}
The argument uses the Fejér identity to connect trigonometric quotients with divisibility and evaluates at integers via the regularized cosine polynomial, while real inputs are handled by the quotient-of-sines representation. A consequential feature is that the Fejér-regularized quotient \(\bigl(\sin(\pi x)/\sin(\pi x/i)\bigr)^2\) yields, after normalization by \(i^2\), an exact divisor filter at integer arguments, \(\phi_i(n)=\mathbf{1}_{\,i\mid n}\). Superposition with weights then turns the filter into divisor sums—most notably \(\sum_i i^{2}\phi_i(n)=\sigma_{2}(n)\)—which provides the structural step enabling the smooth analogues \(\mathcal{P}_{\tau}\) and \(\mathcal{P}_{\sigma}\) via \(W(i)=i^{-2}\) and \(W(i)=i^{-1}\), with limits \(\mathcal{P}_{\tau}(n;\kappa)\to\tau(n)-2\) and \(\mathcal{P}_{\sigma}(n;\kappa)\to\sigma(n)-n-1\).

\subsection{Summary of Results and Open Questions}
Beyond a single indicator, smooth analogues are developed for arithmetic functions such as $\tau(x)$ and $\sigma(x)$. The function $\mathcal{P}_{\tau}$ is $C^\infty$ and approximates integer-prime zeros in the steep-cutoff limit. However, for finite steepness, it exhibits real zeros ("companion zeros") in the neighborhood of primes rather than at the primes themselves. The precise control and potential elimination of these analytic artefacts remains a key open question. For $\mathcal{P}_{\sigma}$, the indicator role is restricted to integer input, as the existence of non-integer zeros cannot be ruled out. The zero of $\mathcal P_\sigma$ to the right of an odd prime $p\ge5$ is shown to be unique for $\kappa\ge3.25\,(p+1)\ln p$ by a computer-assisted argument (Proposition~\ref{prop:Psigma-unique}; interval-arithmetic soundness assumed, and for $p\ge23$ two auxiliary estimates are verified numerically but not written out, see Remark~\ref{rem:right-zero-scope}); its behaviour as $\kappa\to\infty$ and its uniformity in $p$ remain open (Conjecture~\ref{conj:Psigma-companions}).

\subsection{Generalization to the Fej\'er--Dirichlet Lift}
The divisor-filter property, $F(n,i)/i^2 = \mathbf{1}_{i\mid n}$, is the central mechanism of this work. A subsequent manuscript~\cite{fuchs2025-fejer} generalizes this principle by constructing entire functions $\mathcal{T}_a(z) = \sum_{i\ge1} a(i) F(z,i)/i^2$ for arbitrary weight sequences $a(i)$. This operator, termed the \href{https://arxiv.org/abs/2509.12297}{Fej\'er--Dirichlet lift}, provides an analytic interpolant for the Dirichlet convolution $(a*1)(n)$. Its associated Dirichlet series is shown to factor as $\zeta(s)A(s)$, where $A(s)$ is the Dirichlet series of the weights. This framework reveals that the constructions herein are specific instances of a broader structural connection between Fejér analysis and the spectral theory of L-functions.

\subsection{Directions for Future Research}
Further directions include comparisons with other smooth prime indicators and alternative choices of $W(i)$ to model additional arithmetic functions, e.g., $\sigma_k(n)$ via $W(i)=i^{k-2}$. The properties of an analogous indicator for perfect numbers, such as $\mathcal{S}_\sigma(x;\kappa)-(2x-1)$, are left for future investigation. Recursive constructions based on indicator outputs (e.g., for $\omega(n)$) are another option; establishing their properties appears nontrivial. Practical efficiency for large inputs is not the focus here.

